% Round 15. Arbitrary actual top inputs; a curve-specific output theorem.
% This does not assert an auxiliary-factor or arbitrary-base-change formula.
\section{Projection of arbitrary actual top forms to a curve}
\label{r15tc:section}

In relative dimension one over a curve, an actual $(2,2)$ input
is a top form.  This permits the missing vertex contributions to
be retained as ordinary measures.  The argument does not extend
to a higher-dimensional base merely by adding an auxiliary factor.

\begin{lemma}[Equal norm tuples give equal actual presentations]
\label{r15tc:norm-tuples}
Let $V$ be an ordinary analytic open and let $f=(f_1,\ldots,f_N)$
and $r=(r_1,\ldots,r_N)$ be tuples of analytic units there.
If $|f_i|=|r_i|$ pointwise on $V$ for every $i$, then
\[
                 t_f^\star\gamma=t_r^\star\gamma
\tag{TC15.1}\label{r15tc:equal-norm}
\]
for every ambient $\delta$-form $\gamma$ for which these
presentations are defined.  No PWS restriction on $\gamma$
is needed.
\end{lemma}

\begin{proof}
Use the defining tests on every common refined skeleton with
coordinates containing $f$, $r$, and an arbitrary further unit
tuple.  The components of $t_f$ and $t_r$ are literally the same
linear functions on that tropical space, so their tropical
pullbacks of $\gamma$ are equal.  Mihatsch Definition~4.2
identifies the resulting actual forms.  The statement is about
equality of the functions on the whole ordinary open, not merely
equality of their values at one point.
\end{proof}

\begin{lemma}[Realizing a finite graph measure, including atoms]
\label{r15tc:graph-measure}
Let $C/F$ be a smooth geometrically integral projective curve,
viewed in the native $K$ theory for finite $F/K$.  Let $\Gamma$
be a finite connected compact split semistable skeleton of the
kind in frozen Lemma \texttt{r14x:graph-extension}, allowing
finite subdivisions at arbitrary real positions.  Suppose
$\mu$ is a finite signed measure supported on $\Gamma$, whose
edge densities are PWS and smooth up to endpoints on a finite
subdivision, and whose other part is a finite sum of atoms.
Then $\mu$ is the measure of an ordinary actual
$(1,1)$ form on $C$.

This hypothesis does not assume in advance that $\mu$ is the
measure of an actual form.
\end{lemma}

\begin{proof}
Enlarge the model skeleton if necessary so it contains the
support of an ample model curvature $h$; extend the given
measure by zero on the added edges.  The added graph remains
finite and connected.  Put $M_h=\int_C h>0$ and
\[
             \mu_0=\mu-\frac{\mu(\Gamma)}{M_h}\,[h].
\tag{TC15.2}\label{r15tc:mass-normalization}
\]
Its total mass is zero.  The frozen weighted-graph Green lemma
\texttt{r14v:graph-green} applies directly to arbitrary signed
graph-measure data: it integrates each edge density twice,
corrects the endpoint values, and solves the finite weighted
graph Laplacian for the remaining vertex masses.  It therefore
supplies a continuous PWS graph function $f$ with
$D_\Gamma f=\mu_0$.  The proof of that lemma is a graph
calculation and has no actual-representability premise.

Frozen Lemma \texttt{r14x:graph-extension} realizes the
retraction-constant extension of $f$ as an ordinary actual
PWS scalar on $C$.  Its arbitrary real breakpoints include
atoms at type-III points; no assertion that they become model
vertices after a finite extension is needed.  Scalar
Poincar\'e--Lelong gives exactly the native weighted second
derivatives on edges and outgoing slopes at vertices.
Consequently
\[
                \frac{\mu(\Gamma)}{M_h}\,h+Df
\tag{TC15.3}\label{r15tc:realization}
\]
is an ordinary actual form with measure $\mu$.
All residue and edge weights are those of the native graph;
they are not normalized to geometric sheet counts.
\end{proof}

\begin{theorem}[Actual top-form projection over a proper curve]
\label{r15tc:projection}
Let $C/K$ be smooth integral projective of dimension one and
let $\pi:C\times\mathbf P^1\to C$ be projection.  For every
\[
                   \alpha\in B^{2,2}(C\times\mathbf P^1)
\]
there is a unique ordinary actual form
$\mathcal T\alpha\in B^{1,1}(C)$ whose native top measure is
the proper push of the measure of $\alpha$.  In particular
\[
     \int_C\mathcal T\alpha\,\eta
       =\int_{C\times\mathbf P^1}\alpha\wedge\pi^*\eta
\tag{TC15.4}\label{r15tc:pairing}
\]
for every ordinary actual scalar $\eta\in B^{0,0}(C)$,
and hence for all ordinary smooth compact scalar tests.
The construction is linear.  The input need not be PWS.

The theorem concerns a top form over a one-dimensional base.
It does not assert compatibility after adjoining an arbitrary
auxiliary factor or algebraic base change to a higher-dimensional
base.
\end{theorem}

\begin{proof}
\emph{Finite rational presentations.}
Choose finitely many algebraic rational/Laurent affinoids
whose interiors cover the projective surface and ordinary
smooth partition functions compactly supported in those
interiors.  Choose finite unit-moment actual presentations
for $\alpha$ there, and include the moment units of the
partition functions and of the domain inequalities.
As in the frozen curve-output proof, approximate each analytic
unit $f$ in its affinoid algebra by a rational function $R$
with $|f-R|<\min|f|$.  Then $|f|=|R|$ on the whole affinoid
and the latter is still a unit.  Lemma~\ref{r15tc:norm-tuples}
makes this an exact replacement of the actual presentation,
including its positive-delta parts.  The initial choice of
algebraic rational/Laurent domains is essential to the density
of these rational functions.

After a finite native cover $a:C'\to C$, the finitely many
rational functions in the fiber variable split as
\[
                   R_i(x,z)=c_i(x)\prod_j(z-a_j(x))^{n_{ij}}
\tag{TC15.5}\label{r15tc:split}
\]
away from a finite set of base closed points.  Add the marks
$0$ and $\infty$ to the fixed root list, and combine equal
marks.  Include the zeros and poles of the coefficients,
the marks, their nonzero pairwise differences, and all required
domain data in that finite bad set.  The normal projective
curve with the splitting function field is smooth in
characteristic zero.  Enlarge its constant field further so
that the finite marked base curve has a regular split
semistable model.  The full degree of the resulting map $a$
is used below.  Put $Y'=C'\times\mathbf P^1$ and denote
the pulled-back actual top form by $\alpha'$.

\emph{A finite graph supports the pushed measure.}
Mihatsch Corollary~4.9(1), or the frozen ordinary
immersion-germ theorem, makes $\alpha'$ vanish on ordinary
neighborhoods of every bad vertical fiber and every root
graph section, explicitly including the section at infinity:
all are algebraic curves in a surface, whereas
the form has bidegree $(2,2)$.  Compactness gives simultaneous
zero neighborhoods for this finite list.  Choose a marked
semistable base skeleton carrying all base functions in
\eqref{r15tc:split}, and truncate its rays in the zero
neighborhoods over the bad base points.  Properness of the
projection allows a whole base neighborhood of each bad
point to be chosen with zero inverse-image form.  As in the
frozen proper-curve construction, a further finite ramified
extension and regularization subdivide the finitely many
model edges sufficiently to make the model graph simplicial.
Include that extension in $a$ and its native degree.  Let
$\Gamma$ be the resulting compact simplicial model graph,
with any additional real breakpoints treated as subdivisions.

On every complementary ball away from those terminal zero
balls, all coefficient and root-difference norms are constant.
The joint norm image of the functions $z-a_j(x)$ on its
projective-line fibers is the same fixed marked-tree image,
independent of $x$.  Each source moment tuple on its original
presentation domain therefore has image of dimension at most
one over this ball.  This is the point at which a PWS-only
argument would be insufficient.  Instead use Mihatsch
Corollary~4.9(2): an actual top form presented by a tuple $f$
is supported on its full-dimensional skeleton $\Sigma(-,f)$.
That two-dimensional skeleton is empty when the tuple's
moment image has dimension at most one.  The actual top form
is consequently zero on the whole ordinary open in question,
including every positive-delta component.  The same conclusion
on a virtual ball follows after its individual splitting
extension, with faithful pullback; no single field is assumed
to split all residue directions.

Thus, for the native signed Radon measure $\mu'$ of
$\alpha'$, one has
\[
                \operatorname{Supp}|\mu'|
                       \subseteq(\pi')^{-1}(\Gamma).
\tag{TC15.6}\label{r15tc:graph-support}
\]

\emph{Retaining all polyhedral strata.}
Use the simplicial model-PL graph embedding
$j:\Gamma\hookrightarrow\mathbf R^N$ from the frozen
proper-curve graph-extension lemma.  Its coordinates are
global actual PL scalars constant on retraction fibers.
Append their local unit-moment presentations after pulling
them to the source, and take a common finite subdivision
on each compact partition support.  Actual top-form
integration and Corollary~4.9(2) now express that source
measure as finitely many native weighted polyhedral
measures on cells of dimensions zero, one, and two.
The map to $j(\Gamma)$ is piecewise affine on this
subdivision.  We push these measures as measures, retaining
every dimension.  No positive-delta form is pulled through
a merely piecewise affine map.

On a cell whose map has rank one, its pushed measure has
a PWS edge density.  For a one-dimensional cell this is
ordinary affine substitution with its nonzero constant
Jacobian.  For a two-dimensional cell it is a finite sum
of integrals of smooth coefficients over intervals with
affine endpoints, exactly as in the verified moving-interval
calculation.  All coefficients are smooth up to endpoints
after finite subdivision.  A rank-zero cell contributes a
single atom, whose mass is the integral of its coefficient
over that compact cell; a zero-dimensional cell also
contributes an atom.  There are finitely many such terms.
In particular the vertex-supported contribution exhibited
in Proposition~\ref{r15vg:vertex} is retained here.

These assertions use compactly supported partition pieces
inside the presentation domains.  Analytic boundary images
do not cut a nonzero paired measure, and all native
multiplicities are carried by the cell weights.  Refining a
cell creates no new measure; its faces are counted according
to the original trigraded polyhedral measure, not discarded
as if every component were an area density.

For every continuous test function on $C'$, its restriction
to $\Gamma$ factors through the embedding $j$.  Hence this
calculation identifies the genuine Radon push measure
$(\pi')_*\mu'$, rather than only its values on selected
toric tests.  It is supported on $\Gamma$ and consists of
PWS edge densities and finitely many atoms.
Lemma~\ref{r15tc:graph-measure} supplies an ordinary actual
form $\beta'\in B^{1,1}(C')$ with this measure.

\emph{Native finite descent.}
Set $e=\deg_K(a)$, the full native finite degree, and put
\[
               \beta=e^{-1}\operatorname{Tr}^{\mathrm{act}}_a
                                      (\beta').
\tag{TC15.7}\label{r15tc:trace}
\]
This uses only the frozen finite actual trace in degree
$(1,1)$, Theorem \texttt{r14h:strip-trace}.  For every
ordinary smooth scalar test $\eta$ on $C$, its all-actual
test pairing and the native finite degree formula for
$a\times\mathrm{id}$ give
\[
 \int_C\beta\eta
 =e^{-1}\int_{C'}\beta'a^*\eta
 =e^{-1}\int_{Y'}\alpha'\wedge(\pi')^*a^*\eta
 =\int_{C\times\mathbf P^1}\alpha\wedge\pi^*\eta.
\tag{TC15.8}\label{r15tc:native-pairing}
\]
These top forms define Radon measures.  Equality on ordinary
smooth scalar tests identifies the measures, by the same
scalar density used in the frozen curve-top comparison.
It therefore gives \eqref{r15tc:pairing} for every actual
scalar, which is a continuous PWS function.  Alternatively
the finite all-actual trace and degree formulas apply to
that scalar directly, since the intermediate identity is
already an equality of Radon measures.  Finally actual
current injectivity proves uniqueness, linearity, and
independence of all choices.
\end{proof}

\begin{corollary}[A projective-line bundle over a curve]
\label{r15tc:bundle}
Theorem~\ref{r15tc:projection} also holds for
$\pi:\mathbf P_C(N)\to C$, where $N$ is a rank-two
algebraic vector bundle and the input is any global actual
top $(2,2)$ form.
\end{corollary}

\begin{proof}
Trivialize $N$ on the complement of finitely many base
closed points.  The actual top form vanishes on an ordinary
neighborhood of every omitted fiber, by its degree and the
ordinary immersion-germ theorem.  Properness shrinks these
to base neighborhoods with zero entire inverse-image form.
Transport it through the algebraic trivialization and
extend by zero to $C\times\mathbf P^1$.  This is a literal
ordinary actual form, since it is already zero before any
omitted fiber is reached.  Its projected measure is the
original projected measure.  Apply the theorem.
\end{proof}

\paragraph{Dependencies and scope.}
The two primary points used for the positive-delta extension
are Mihatsch Definition~4.2 and Corollary~4.9(2), in
\url{https://arxiv.org/html/2112.10018v4}.
The finite graph geometry and rational factorizations are
those of
\path{round14/projective_cohomology/relative_curve_top_output_and_comparison.tex}.
The realization is the explicit graph-measure lemma
\texttt{r14v:graph-green} in
\path{round14/structured_descent/projective_line_aeppli_and_surface_descent.tex},
followed by \texttt{r14x:graph-extension} in
\path{round14/structured_descent/proper_curve_aeppli.tex}.
It does not invoke the conclusion of the old curve-potential
theorem on a measure whose actual nature has not yet been
proved.

The result independently resolves arbitrary actual $(2,2)$
projection over a proper curve, with every actual scalar test.
In higher base dimension the input is not top degree, so this
particular supporting-skeleton proof does not apply.  The later
Theorem~\ref{r15a2:projection} resolves that higher-base output
by the universal kernel and finite ambient atlas, with all
ordinary smooth compact mixed tests.  It is a separate argument.
The all-actual scalar comparison proved here does not become
an all-actual proper pushforward theorem in general, and no
auxiliary-factor mixed comparison is inferred from this
curve-top proof alone.
